% Corrected version of the "origin of the energy surplus" block. Changes relative to
% your draft: typos (suppress, magnitude, energy, principal), grammar (tells, the more
% opportunities), the trace figure swapped to the eta=1 real sampler version so the
% Re=2000 overshoot claim holds, the embedding regimes corrected to the four basis
% regimes, PC1 variance filled in at 96 percent, the numbered markers described as the
% first clean estimate of each chain rather than its end, and the subfigure captions
% filled. Sentence structure kept as close to yours as the fixes allow.

\paragraph{Uncovering the origin of this energy surplus at inference time} In order to
suppress the fine-scale energy surplus provided by turbulence-promoting models, we need
to understand the inference time mechanisms promoting the injection of these
turbulences. Figure \ref{fig:chain-energy-trace-K3} depicts how models tuned in
different regimes control the amount of $k\in[32,96)$ band energy they inject into
their respective generations through the three successive denoising chains they follow,
given the same sample from the $Re=1000$ pool and the production sampler. The magnitude
of the fine-scale energy levels after the first denoising step tells us that the more
turbulent the model's home regime is, the more fine-scale energy it injects by
leveraging larger magnitude noise predictions $\epsilon_\theta$. All models follow the
same energy injection pattern, raising the energy content of the sample at first,
before flattening and slightly reversing as the chain refines towards a clean sample.
After the first denoising step, the $Re=4000$ and $Re=8000$ finetunes already overshoot
the ground truth fine-scale content of the sample, while the $Re=2000$ falls just short
of it at the end of its first chain and the $Re=1000$ finetune remains farther
below.\newline
The re-noising events of the frozen inference configuration then allow each model to
inject further fine-scale energy into the flow, pushing the $Re=2000$ model past the
truth over the remaining chains. Importantly, we note that the larger the noising
increment $t$, the larger the magnitude of the fine-scale energy injection is.
Equally, the more re-noising events there are, the more opportunities the model has to
push the reconstruction towards the energy levels of its home regime.\newline

To support this hypothesis we embed per sample log energy spectra of the regimes
$Re\in\{1000, 2000, 4000, 8000\}$ in a compact space using principal component
analysis. The first principal component explains $96\%$ of the variance in the dataset
and roughly aligns with the flow regime. In Figure \ref{fig:cross-manifold-K3}, we
project the trajectories of the models tested onto this manifold, marking the first
clean estimate of each denoising chain and the final reconstruction. First, we note
that the base model, which is unable to inject fine-scale energy in the flow, generates
samples which stand close together and near the DDIM recon over the manifold. The
$Re=2000$ model reaches the $Re=1000$ cloud within its first chain while closely
matching its fine-scale energy. By the opening of its second chain, its generation has
traveled in the positive direction of PC1, closer to where $Re=2000$ samples lie on the
data manifold. The $Re=1000$ finetune first generates a sample close to the base
reconstruction, before drifting along PC1 as more fine-scale energy is injected into
it. The $Re=4000$ and $Re=8000$ finetunes over-dose fine-scale energy upon their first
chain, and drift farther along the PC1 axis, each getting closer to where its home
regime samples lie over the manifold. This figure illustrates how specialist models
force the generated samples to drift towards their home regime over the data manifold,
and the jump between the end of one chain and the opening of the next is the injection
itself. Similarly to the observation made over Figure \ref{fig:chain-energy-trace-K3},
the larger the $t$, the farther the clean sample can travel from its previous clean
realisation.\newline

\begin{figure}[H]
    \centering
    \begin{subfigure}[b]{\textwidth}
        \centering
        \includegraphics[width=\textwidth]{Figures/Cross-regime/chain_energy_trace_eta1.pdf}
        \caption{Fine band energy of the running clean estimate through the three
        chains, real sampler, identical noise across models.}
        \label{fig:chain-energy-trace-K3}
    \end{subfigure}
    \\[1ex]
    \begin{subfigure}[b]{\textwidth}
        \centering
        \includegraphics[width=\textwidth]{Figures/Cross-regime/offmanifold_topview_re1000.pdf}
        \caption{The same chains on the spectral manifold, target $Re=1000$. Numbered
        dots mark the first clean estimate of each chain, diamonds the final
        reconstruction.}
        \label{fig:cross-manifold-K3}
    \end{subfigure}
    \caption{The injection of fine scale energy through the frozen K3 configuration,
    seen through the band energy of the clean estimate and through the corresponding
    movement over the data manifold.}
    \label{fig:cross-manifold-move}
\end{figure}

We then leverage different test time inference configurations to promote downward
transfer of more energetic models, and observe the repercussions over the energy
spectrum and the trajectories over the data manifold.
